% !PRINTED 403
\weberpage{403}
%<<CONT
\[
\frac{1}{1}, \frac{2}{1}, \frac{7}{4}, \frac{16}{9}, \frac{23}{13}, \frac{2201}{1244}, \frac{4425}{2501} \cdots
\]
Wenn wir die beiden letzten Brüche in Decimalbrüche verwandeln, so erhalten wir für $x$ die beiden Grenzen:
\[
1{,}7692926; \quad 1{,}76929228,
\]
und der Fehler in der letzteren etwas zu kleinen Zahl ist kleiner als
\[
0{,}00000016.
\]
Dass man hier nach wenig Schritten ein gutes Resultat erhält, beruht auf dem Umstande, dass hier ziemlich früh ein grösserer Theilnenner 95 auftritt. In den meisten Fällen bekommt man für die Theilnenner kleine Zahlen; dann wachsen die Nenner der Näherungsbrüche langsam, und man erhält nur mühselig ein Resultat von grösserer Genauigkeit.

Die Auffindung grösserer Theilnenner, z. B. 95, wird durch die Bemerkung erleichtert, dass, wenn eine Wurzel einer Gleichung einen grossen Werth hat, der negative Co\"efficient der zweithöchsten Potenz der Unbekannten eine, wenn auch nur rohe Annäherung an diese Wurzel giebt, wie z. B. in der vorletzten der obigen Gleichungen 94 an 95.

\begin{center}
\S.\,130.\\[0.3em]
\gesperrt{Rationale Wurzeln ganzzahliger Gleichungen.}\\
\gesperrt{Reducible Gleichungen.}
\end{center}

Wir schliessen diesen Abschnitt mit einer Betrachtung über ganzzahlige Gleichungen, die, wenn sie auch nicht unmittelbar mit der Theorie der Kettenbrüche in Beziehung steht, doch hier am passendsten eine Stelle findet.

Wenn man es mit Gleichungen zu thun hat, deren Co\"efficienten rationale Zahlen sind, so wird man zunächst nach etwaigen rationalen Wurzeln suchen. Wenn die Gleichung
\[
f(x) = a_0 x^n + a_1 x^{n-1} + a_2 x^{n-2} + \cdots + a_{n-1} x + a_n = 0
\]
rationale Co\"efficienten hat, so können wir immer annehmen, dass diese Co\"efficienten ganze Zahlen sind; man hat nur n\"othig, um den Fall gebrochener Co\"efficienten darauf zurückzuführen, die ganze Gleichung mit einem gemeinschaftlichen Vielfachen aller Nenner zu multipliciren.
